% ====================================================================
% BAB 4: HASIL DAN PEMBAHASAN
% ====================================================================
% File ini berisi hasil penelitian dan pembahasan yang mencakup:
% - Hasil penelitian (data terolah dan temuan utama)
% - Pembahasan hasil (interpretasi dan analisis kritis)
%
% EDIT konten sesuai dengan hasil penelitian Anda.
% ====================================================================

\chapter{HASIL DAN PEMBAHASAN}


\section{Pengumpulan Data}
\label{sec:data-collection}

Data penelitian diperoleh dari grup percakapan WhatsApp [nama/deskripsi grup] setelah memperoleh izin dari admin grup sebagaimana dijelaskan pada Lampiran~\ref{subsec}. 
Pengumpulan data dilakukan menggunakan fitur \textit{Export Chat} bawaan aplikasi WhatsApp, yang menghasilkan berkas dalam format \textit{.zip} berisi riwayat percakapan dalam satu berkas \textit{.txt} tanpa media. 
Berkas \textit{.zip} tersebut terlebih dahulu diekstraksi untuk memperoleh file \textit{.txt} yang selanjutnya digunakan sebagai sumber data utama penelitian. 
Data yang diekspor mencakup periode percakapan dari [tanggal awal] hingga [tanggal akhir], dengan total [jumlah] pesan dari [jumlah] partisipan grup.

Sebagai ilustrasi, potongan data mentah hasil ekspor ditampilkan pada Listing~\ref{lst:contoh-data-mentah}. 
Secara umum, setiap baris percakapan mengikuti pola \texttt{DD/MM/YY HH.MM - Pengirim: Pesan}. 
Namun demikian, data mentah juga memuat variasi format lain seperti pesan sistem (misalnya anggota ditambahkan ke grup), placeholder media, tautan, dan penanda pesan yang telah diedit.

\begin{lstlisting}[caption={Contoh potongan data mentah hasil ekspor WhatsApp (nama disamarkan)}, label={lst:contoh-data-mentah}, basicstyle=\ttfamily\small]
02/02/26 20.46 - User-1: Woi Katakan pada dunia hsr back
02/02/26 20.47 - User-1: <Media tidak disertakan>
27/09/25 10.35 - User-3: sheesh
03/02/26 11.24 - User-4: https://link
28/09/25 09.48 - User-5: gara gara dapet evernight jadi pengen hyacine <Pesan ini diedit>
24/04/26 05.12 - User-7 sudah ditambahkan
\end{lstlisting}

Berkas \textit{.txt} selanjutnya diproses melalui tahap \textit{parsing} menggunakan \textit{regular expression} untuk mengekstrak tiga atribut utama, yaitu \textit{Timestamp}, \textit{Pengirim}, dan \textit{Pesan}. 
Pada tahap ini, nilai \textit{timestamp} dinormalisasi ke format ISO 8601 (\texttt{YYYY-MM-DDTHH:MM}) untuk memudahkan pengurutan kronologis, filtering temporal, dan analisis berbasis waktu.

Sebagai ilustrasi, hasil \textit{parsing} dari potongan data mentah pada Listing~\ref{lst:contoh-data-mentah} disajikan pada Tabel~\ref{tab:contoh-hasil-parsing}.

\begingroup
\pagestyle{fancy}
\renewcommand{\arraystretch}{1}
\begin{table}[H]
\centering
\caption{Contoh Hasil Parsing Data WhatsApp ke Format CSV}
\label{tab:contoh-hasil-parsing}
\scriptsize
\begin{tabular}{p{12cm}}
\toprule
\textbf{Pesan} \\
\midrule
Woi Katakan pada dunia hsr back \\
sheesh \\
https://link \\
gara gara dapet evernight jadi pengen hyacine \\
\bottomrule
\end{tabular}
\end{table}
\endgroup

\section{Praproses Data}
\label{sec:preprocessing}
\noindent Uraikan langkah pra-pemrosesan yang dilakukan sesuai metodologi:
\begin{enumerate}[nosep]
    \item Pemfilteran pesan sistem dan pesan media
    \item Normalisasi (lower-case, mention → @USER, URL → HTTP)
    \item Transliteration emotikon dan normalisasi slang
    \item Tokenisasi dan stopword removal (Sastrawi + manual)
\end{enumerate}
Berikan contoh before/after satu atau dua pesan untuk ilustrasi.

\section{Pembuatan Model}
\label{sec:descriptive-results}
\noindent Sajikan statistik deskriptif utama setelah pra-pemrosesan: mean, median, standar deviasi panjang pesan, frekuensi kata teratas, dan distribusi topik awal jika relevan.

\subsection{Embedding dan Reduksi Dimensi}
\label{subsec:embedding-umap}
\noindent Jelaskan penggunaan IndoBERTweet untuk pembuatan embedding, parameter penting yang dipilih, dan hasil reduksi dimensi menggunakan UMAP. Sertakan visualisasi (scatter plot UMAP) dan interpretasi singkat.

\subsection{Clustering dan Representasi Topik}
\label{subsec:clustering-topics}
\noindent Paparkan penggunaan BIRCH untuk clustering dan BM25 untuk representasi topik. Untuk setiap topik utama, tampilkan:
\begin{itemize}[nosep]
    \item Kata kunci representatif (top-n)
    \item Jumlah dokumen per klaster
    \item Contoh pesan yang termasuk dalam topik
\end{itemize}

\subsection{Evaluasi Kuantitatif Model}
\label{subsec:model-evaluation}
\noindent Tampilkan hasil metrik evaluasi yang digunakan: topic coherence, topic diversity, embedding density, dan intra-topic similarity. Sajikan hasil dalam tabel per skenario parameter (baseline vs modifikasi) dan beri interpretasi singkat.

\subsection{Perbandingan Eksperimen (Baseline vs Modifikasi)}
\label{subsec:experiments-comparison}
\noindent Bandingkan kinerja skenario berbeda (mis. HDBSCAN vs BIRCH, c-TF-IDF vs BM25, model embedding berbeda). Sertakan tabel ringkasan dan highlight perbedaan signifikan.

\subsection{Uji Asumsi dan Analisis Inferensial}
\label{subsec:inferential-results}
\noindent Jika relevan, laporkan uji normalitas, uji homogenitas, dan hasil pengujian hipotesis. Jelaskan metode statistik yang digunakan dan implikasinya terhadap interpretasi hasil.

\subsection{Implementasi Sistem dan Evaluasi Fungsional}
\label{subsec:system-evaluation}
\noindent Laporkan hasil pengujian aplikasi berbasis web: keandalan upload, waktu proses end-to-end, pengalaman pengguna (jika ada pengujian usability), serta contoh tampilan hasil.

\subsection{Analisis Hasil}
\label{subsec:result-analysis}
\noindent Analisis sintesis hasil kuantitatif dan kualitatif: apakah topik yang dihasilkan bermakna, stabil, dan sesuai konteks grup WhatsApp. Kaitkan dengan tujuan penelitian.

% --- BAGIAN 4.2: PEMBAHASAN ---
\section{Pembahasan}
\label{sec:discussion}

\subsection{Interpretasi Temuan Utama}
\label{subsec:interpretation}
\noindent Paparkan penjelasan mendalam mengenai temuan utama, mengapa hasil tersebut muncul, dan bagaimana temuan ini menjawab rumusan masalah. Hubungkan temuan dengan teori dan tinjauan pustaka di Bab 2.

\subsection{Analisis Kualitatif Topik}
\label{subsec:qualitative-topic-analysis}
\noindent Bahas beberapa topik penting secara kualitatif: konteks sosialnya, contoh kutipan pesan, dan relevansi bagi pengguna.

\subsection{Implikasi Praktis dan Rekomendasi Implementasi}
\label{subsec:practical-implications}
\noindent Jelaskan implikasi untuk pengembangan sistem dan pengguna akhir, termasuk rekomendasi konfigurasi model, antarmuka, dan fitur tambahan.

\subsection{Keterbatasan Penelitian}
\label{subsec:limitations}
\noindent Sebutkan keterbatasan data, metode, dan eksperimen (mis. satu grup WhatsApp, potensi bias, keterbatasan anotasi) serta dampaknya pada generalisasi hasil.

\subsection{Saran dan Pekerjaan Lanjutan}
\label{subsec:future-work}
\noindent Rekomendasi teknis dan riset lanjutan: perluasan dataset, eksperimen model embedding lain, perbaikan pipeline pra-pemrosesan, atau pengujian usability yang lebih formal.

\subsection{Ringkasan Hasil Utama}
\label{subsec:summary-results}
\noindent Ringkas temuan paling penting dalam 3--5 poin yang langsung menjawab rumusan masalah.

% ====================================================================
% CATATAN: Isi tiap subsection diisi dengan hasil eksperimen dan pembahasan spesifik Anda.
% Gunakan tabel dan gambar sesuai kebutuhan, dan rujuk ke Bab 3 untuk metodologi yang relevan.
% ====================================================================

% ====================================================================
% CATATAN PENGISIAN BAB HASIL & PEMBAHASAN:
%
% 1. HASIL PENELITIAN:
%    - Presentasikan data secara objektif (tidak interpretasi)
%    - Gunakan tabel dan grafik untuk clarity
%    - Jelaskan setiap temuan utama
%    - Gunakan statistik deskriptif dengan tepat
%
% 2. PEMBAHASAN - SANGAT PENTING:
%    - JANGAN hanya repeat hasil, JELASKAN MENGAPA
%    - Hubungkan dengan teori dan literature dari Bab 2
%    - Analisis kritis: apakah sejalan/berbeda dengan penelitian terdahulu?
%    - Jelaskan implikasi praktis
%    - Sertakan keterbatasan penelitian
%    - Inilah bagian "value-add" dari skripsi Anda
%
% 3. FORMAT:
%    - Gunakan subsection untuk mengorganisir topik
%    - Gunakan referensi ke tabel/gambar: "seperti terlihat pada Tabel \ref{tab:label}"
%    - Sertakan kutipan dari literature: \citet{Author2020}
%
% 4. ELEMEN KRITIS:
%    - Analisis mendalam (bukan permukaan)
%    - Kritisisme terhadap hasil: apa artinya? kenapa terjadi?
%    - Perbandingan dengan studi sebelumnya
%    - Keterbatasan yang jujur dan transparan
% ====================================================================
